
We investigated gradient-based residual-corrected Jacobian stable rank regularization during GPT-2 pretraining. The implementation combined Hutchinson trace estimation (10 probes) and power iteration (5 iterations) differentiated through PyTorch autodiff. Results show comprehensive failure: zero stable rank measurements, perplexity increase from 59.3 to 45,792, and regularization losses of -17.5 billion. Baseline training succeeded under identical conditions, isolating the failure to the regularization mechanism.

Root cause analysis identified insufficient Hutchinson probes for 768-dimensional embeddings, power iteration non-convergence in residual-corrected graphs, and gradient pathologies in deep autodiff chains. The measurement-control coupling means noisy estimates directly corrupt training signals, creating a feedback loop that prevents convergence.

This negative result documents specific failure modes in gradient-based spectral control and suggests post-hoc SVD-based observational studies as a numerically stable alternative for testing whether Jacobian stable rank correlates with efficiency properties.
